\textbf{Phần 1. Trắc nghiệm lựa chọn}
\setcounterpageref{ex}{0}

\begin{ex} %[5X11N67]
Biểu đồ hình quạt tròn biểu diễn toàn bộ tổng thể tương ứng với bao nhiêu phần trăm?
\choice
{\True $100\%$}
{$50\%$}
{$200\%$}
{$360\%$}
\loigiai{
Toàn bộ hình quạt tròn biểu diễn $100\%$.\\
\textbf{Đáp án của câu hỏi là A}
}
\end{ex}

\begin{ex} %[5X11N67]
Gieo một con xúc xắc $30$ lần, mặt $5$ chấm xuất hiện $6$ lần. Tỉ số của số lần xuất hiện mặt $5$ chấm là
\choice
{$\dfrac{5}{30}$}
{$\dfrac{1}{6}$}
{\True $\dfrac{1}{5}$}
{$\dfrac{2}{5}$}
\loigiai{
$\dfrac{6}{30} = \dfrac{1}{5}$.\\
\textbf{Đáp án của câu hỏi là C}
}
\end{ex}

\begin{ex} %[5X11H67]
Biểu đồ hình quạt tròn cho biết cơ cấu phương tiện đến trường của $400$ học sinh: Xe đạp $50\%$, Đi bộ $30\%$, Xe buýt $20\%$. Số học sinh đi xe buýt là
\choice
{$60$ bạn}
{\True $80$ bạn}
{$100$ bạn}
{$120$ bạn}
\loigiai{
$400 \times 20 : 100 = 80$ (bạn).\\
\textbf{Đáp án của câu hỏi là B}
}
\end{ex}

\begin{ex} %[5X11H67]
Tung một đồng xu $50$ lần có $22$ lần xuất hiện mặt Sấp. Tỉ số phần trăm số lần xuất hiện mặt Ngửa là
\choice
{$44\%$}
{$50\%$}
{$48\%$}
{\True $56\%$}
\loigiai{
Mặt Ngửa xuất hiện: $50 - 22 = 28$ lần $\Rightarrow 28 : 50 = 56\%$.\\
\textbf{Đáp án của câu hỏi là D}
}
\end{ex}

\begin{ex} %[5X11H67]
Cho bảng số liệu điểm kiểm tra của một tổ: $9, 8, 7, 10, 8, 9, 8, 9$. Điểm xuất hiện nhiều lần nhất là
\choice
{$7$}
{\True $8$ (và $9$)}
{$9$}
{$10$}
\loigiai{
Điểm 8 và điểm 9 đều xuất hiện 3 lần.\\
\textbf{Đáp án của câu hỏi là B}
}
\end{ex}

\begin{ex} %[5X11H67]
Khảo sát sở thích đọc sách của $150$ học sinh: Truyện tranh $60\%$, Sách khoa học $25\%$, Sách lịch sử $15\%$. Số bạn thích sách lịch sử là
\choice
{$15$ bạn}
{$25$ bạn}
{$30$ bạn}
{\True $22{,}5$ bạn}
\loigiai{
$150 \times 15 : 100 = 22{,}5$ (bạn).\\
\textbf{Đáp án của câu hỏi là D}
}
\end{ex}

\begin{ex} %[5X11V67]
Gieo một con xúc xắc $100$ lần, mặt có số chấm là số nguyên tố ($2, 3, 5$) xuất hiện $52$ lần. Tỉ số phần trăm số lần xuất hiện mặt không phải số nguyên tố là
\choice
{\True $48\%$}
{$52\%$}
{$50\%$}
{$46\%$}
\loigiai{
Số lần không phải số nguyên tố: $100 - 52 = 48$ lần $\Rightarrow 48\%$.\\
\textbf{Đáp án của câu hỏi là A}
}
\end{ex}

\begin{ex} %[5X11V67]
Biểu đồ quạt biểu diễn diện tích trồng trọt: Lúa $55\%$, Ngô $25\%$, Khoai $20\%$. Biết diện tích trồng Lúa là $11\text{ ha}$. Tổng diện tích trồng trọt là
\choice
{$18\text{ ha}$}
{$22\text{ ha}$}
{\True $20\text{ ha}$}
{$25\text{ ha}$}
\loigiai{
Tổng diện tích: $11 : 55 \times 100 = 20\text{ (ha)}$.\\
\textbf{Đáp án của câu hỏi là C}
}
\end{ex}

\begin{ex} %[5X11V67]
Khảo sát số giờ sử dụng điện thoại mỗi ngày của $8$ người: $1\text{h}, 2\text{h}, 1{,}5\text{h}, 3\text{h}, 2{,}5\text{h}, 2\text{h}, 1\text{h}, 3\text{h}$. Số giờ sử dụng trung bình là
\choice
{$1{,}8\text{ giờ}$}
{\True $2\text{ giờ}$}
{$2{,}2\text{ giờ}$}
{$2{,}5\text{ giờ}$}
\loigiai{
Tổng: $1 + 2 + 1{,}5 + 3 + 2{,}5 + 2 + 1 + 3 = 16$ giờ.\\
Trung bình: $16 : 8 = 2$ giờ.\\
\textbf{Đáp án của câu hỏi là B}
}
\end{ex}

\begin{ex} %[5X11V67]
Một hộp có $5$ viên bi xanh, $3$ viên bi đỏ và $2$ viên bi vàng. Bạn Nam lấy ngẫu nhiên $1$ viên bi, ghi màu rồi trả lại hộp, làm như vậy $50$ lần thì có $24$ lần lấy được bi xanh. Tỉ số phần trăm thực nghiệm lấy được bi xanh là
\choice
{$40\%$}
{$50\%$}
{\True $48\%$}
{$52\%$}
\loigiai{
$24 : 50 = 48\%$.\\
\textbf{Đáp án của câu hỏi là C}
}
\end{ex}

\begin{ex} %[5X11V67]
Biểu đồ quạt tròn cho biết tỉ lệ bình chọn tiết mục văn nghệ: Múa $35\%$, Hát $40\%$, Kịch $25\%$. Biết có $80$ phiếu bình chọn cho tiết mục Hát. Tổng số phiếu bình chọn là
\choice
{\True $200$ phiếu}
{$160$ phiếu}
{$240$ phiếu}
{$250$ phiếu}
\loigiai{
Tổng số phiếu: $80 : 40 \times 100 = 200$ (phiếu).\\
\textbf{Đáp án của câu hỏi là A}
}
\end{ex}

\begin{ex} %[5X11C67]
Một cuộc điều tra về số lượng vật nuôi của $50$ hộ gia đình:
- $15$ hộ nuôi chó
- $20$ hộ nuôi mèo
- $5$ hộ nuôi cả chó và mèo
- Số hộ còn lại không nuôi con nào.
Tỉ số phần trăm số hộ gia đình không nuôi con vật nào là
\choice
{$30\%$}
{$35\%$}
{$25\%$}
{\True $40\%$}
\loigiai{
Số hộ có nuôi ít nhất 1 con vật: $15 + 20 - 5 = 30$ hộ.\\
Số hộ không nuôi con vật nào: $50 - 30 = 20$ hộ.\\
Tỉ số phần trăm: $20 : 50 = 40\%$.\\
\textbf{Đáp án của câu hỏi là D}
}
\end{ex}

\textbf{Phần 2. Trắc nghiệm đúng sai}
\setcounterpageref{ex}{0}

\begin{ex} %[5X11H67]
Biểu đồ hình quạt tròn biểu diễn kết quả phân loại rác thải tái chế: Giấy $40\%$, Nhựa $35\%$, Kim loại $15\%$, Thủy tinh $10\%$. Tổng khối lượng thu gom là $200\text{ kg}$. Xét tính đúng sai:
\choiceTF
{\True Khối lượng giấy thu gom được là $80\text{ kg}$.}
{\True Khối lượng nhựa thu gom được là $70\text{ kg}$.}
{Khối lượng kim loại thu gom được là $40\text{ kg}$.}
{\True Khối lượng thủy tinh thu gom được là $20\text{ kg}$.}
\loigiai{
\begin{itemize}
    \item Giấy: $200 \times 40\% = 80\text{ kg}$. Đúng.
    \item Nhựa: $200 \times 35\% = 70\text{ kg}$. Đúng.
    \item Kim loại: $200 \times 15\% = 30\text{ kg} \neq 40\text{ kg}$. Sai.
    \item Thủy tinh: $200 \times 10\% = 20\text{ kg}$. Đúng.
\end{itemize}
\begin{itemize}
    \item a) \textbf{Đúng}.
    \item b) \textbf{Đúng}.
    \item c) \textbf{Sai}.
    \item d) \textbf{Đúng}.
\end{itemize}
}
\end{ex}

\begin{ex} %[5X11V67]
Tung một con xúc xắc $80$ lần. Kết quả: Mặt $1$ chấm ($12$ lần), Mặt $2$ chấm ($14$ lần), Mặt $3$ chấm ($10$ lần), Mặt $4$ chấm ($16$ lần), Mặt $5$ chấm ($18$ lần), Mặt $6$ chấm ($10$ lần). Xét tính đúng sai:
\choiceTF
{\True Số lần xuất hiện mặt chẵn là $40$ lần ($50\%$).}
{\True Số lần xuất hiện mặt lẻ là $40$ lần ($50\%$).}
{Mặt có số lần xuất hiện nhiều nhất là mặt $4$ chấm.}
{\True Mặt có số lần xuất hiện nhiều nhất là mặt $5$ chấm ($18$ lần).}
\loigiai{
\begin{itemize}
    \item Mặt chẵn (2, 4, 6): $14 + 16 + 10 = 40$ lần ($50\%$). Đúng.
    \item Mặt lẻ (1, 3, 5): $12 + 10 + 18 = 40$ lần ($50\%$). Đúng.
    \item Mặt 5 chấm xuất hiện 18 lần, nhiều nhất (không phải 4 chấm). Ý 3 Sai, Ý 4 Đúng.
\end{itemize}
\begin{itemize}
    \item a) \textbf{Đúng}.
    \item b) \textbf{Đúng}.
    \item c) \textbf{Sai}.
    \item d) \textbf{Đúng}.
\end{itemize}
}
\end{ex}

\textbf{Phần 3. Trả lời ngắn}
\setcounterpageref{ex}{0}

\begin{ex} %[5X11H67]
Trên biểu đồ quạt tròn, phần thể thao chiếm $35\%$, nghệ thuật chiếm $25\%$, phần còn lại là học thuật. Hỏi phần học thuật chiếm bao nhiêu phần trăm? (Chỉ ghi số).
\shortans{40}
\loigiai{
\textbf{Đáp số:} $40$\\
$100 - (35 + 25) = 40\%$.
}
\end{ex}

\begin{ex} %[5X11H67]
Gieo xúc xắc $50$ lần, có $9$ lần xuất hiện mặt 1 chấm. Tính tỉ số phần trăm số lần xuất hiện mặt 1 chấm. (Chỉ ghi số).
\shortans{18}
\loigiai{
\textbf{Đáp số:} $18$\\
$9 : 50 = 18\%$.
}
\end{ex}

\begin{ex} %[5X11V67]
Biểu đồ quạt tròn cho biết $30\%$ số học sinh của trường tham gia CLB Tiếng Anh, tương ứng $180$ học sinh. Tính tổng số học sinh của trường.
\shortans{600}
\loigiai{
\textbf{Đáp số:} $600$\\
$180 : 30 \times 100 = 600$.
}
\end{ex}

\begin{ex} %[5X11V67]
Tung đồng xu $80$ lần thấy xuất hiện mặt Sấp $36$ lần. Tính tỉ số phần trăm số lần xuất hiện mặt Ngửa. (Chỉ ghi số).
\shortans{55}
\loigiai{
\textbf{Đáp số:} $55$\\
Mặt Ngửa: $80 - 36 = 44$ lần $\Rightarrow 44 : 80 = 55\%$.
}
\end{ex}

\begin{ex} %[5X11V67]
Điểm kiểm tra của $5$ bạn học sinh là: $8, 9, 7, 10, 6$. Tính điểm trung bình của $5$ bạn.
\shortans{8}
\loigiai{
\textbf{Đáp số:} $8$\\
$(8 + 9 + 7 + 10 + 6) : 5 = 40 : 5 = 8$.
}
\end{ex}

\begin{ex} %[5X11C67]
Biểu đồ quạt tròn cho thấy $45\%$ số cây trong vườn là cam, $35\%$ là xoài, còn lại là $60$ cây bưởi. Tính tổng số cây có trong vườn.
\shortans{300}
\loigiai{
\textbf{Đáp số:} $300$\\
Tỉ lệ bưởi: $100\% - (45\% + 35\%) = 20\%$.\\
Tổng số cây: $60 : 20 \times 100 = 300$ cây.
}
\end{ex}
